\documentclass[11pt]{article}
\usepackage[margin=1in]{geometry}
\usepackage{enumitem}
% =============================================================================
% Preambles/header.tex — shared LaTeX/Beamer preamble for this project
%
% Use from a Beamer lecture by prepending:
%     \input{header}
% and compiling with TEXINPUTS=../Preambles:$TEXINPUTS (the /compile-latex
% skill does this automatically).
%
% PALETTE CONTRACT. Color names defined here MUST match the names in
% Quarto/theme-template.scss so Beamer and Quarto renderings use the same
% palette. The scripts/check-palette-sync.sh script enforces this.
%
% To customize: edit the HEX values below AND the matching $variable in
% Quarto/theme-template.scss, then run scripts/check-palette-sync.sh.
% =============================================================================

% -----------------------------------------------------------------------------
% Engine detection + encoding
%
% Our /compile-latex skill uses XeLaTeX, where inputenc is unnecessary and
% can produce encoding warnings. Only load inputenc under pdfTeX.
% -----------------------------------------------------------------------------
\usepackage{iftex}
\ifPDFTeX
  \usepackage[utf8]{inputenc}
\fi

\usepackage{amsmath, amssymb, amsthm}
\usepackage{booktabs}
\usepackage{graphicx}

% -----------------------------------------------------------------------------
% PALETTE — mirrors Quarto/theme-template.scss variable names.
% Load xcolor and define colors BEFORE anything that references them
% (notably hyperref, hypersetup, and the Beamer color assignments).
% -----------------------------------------------------------------------------
\usepackage{xcolor}

\definecolor{primary-blue}{HTML}{012169}      % headings, accents
\definecolor{primary-gold}{HTML}{B9975B}      % emphasis, borders
\definecolor{highlight-yellow}{HTML}{F2A900}  % markers, alerts
\definecolor{light-bg}{HTML}{E8EDF5}          % title / section backgrounds
\definecolor{jet}{HTML}{1A1A1A}               % body text

% Semantic colors (used in diagrams and annotations)
\definecolor{positive}{HTML}{15803D}          % good / observed / identified
\definecolor{negative}{HTML}{B91C1C}          % bad / problematic / bias
\definecolor{neutral}{HTML}{525252}           % reference / context

% Highlight accents (match .hi-* classes in the SCSS)
\definecolor{hi-slate}{HTML}{314F4F}
\definecolor{hi-green}{HTML}{2E8B57}
\definecolor{hi-red}{HTML}{C41E3A}

% Semantic aliases so skill templates and snippets can write
% \textcolor{accent}{...} without hard-coding "primary-blue".
\colorlet{accent}{primary-blue}
\colorlet{accent2}{primary-gold}

% -----------------------------------------------------------------------------
% Hyperref — loaded AFTER xcolor + palette so link colors resolve.
% -----------------------------------------------------------------------------
\usepackage{hyperref}
\hypersetup{colorlinks=true, linkcolor=primary-blue, urlcolor=primary-blue,
            citecolor=primary-blue, pdfborder={0 0 0}}

% -----------------------------------------------------------------------------
% Course code — set once per deck (after \input{header}, before \begin{document})
% via \coursecode{CS401}. Shown in the footer on every slide so a deck's course
% is identifiable without opening the title page. Empty by default.
% -----------------------------------------------------------------------------
\makeatletter
\newcommand{\coursecode}[1]{\gdef\@coursecodetext{#1}}
\gdef\@coursecodetext{}
\makeatother

% -----------------------------------------------------------------------------
% Beamer theme assignments (only apply under Beamer).
%
% We detect Beamer via \@ifclassloaded{beamer}, which is the documented,
% non-internal way. This must be wrapped in \makeatletter/\makeatother
% because \@ifclassloaded contains an @-catcode character.
% -----------------------------------------------------------------------------
\makeatletter
\@ifclassloaded{beamer}{%
  \usefonttheme{professionalfonts}
  \setbeamercolor{structure}{fg=primary-blue}
  \setbeamercolor{frametitle}{fg=primary-blue}
  \setbeamercolor{title}{fg=primary-blue}
  \setbeamercolor{subtitle}{fg=primary-gold}
  \setbeamercolor{author}{fg=primary-blue}
  \setbeamercolor{institute}{fg=neutral}
  \setbeamercolor{date}{fg=primary-gold}
  \setbeamercolor{itemize item}{fg=primary-blue}
  \setbeamercolor{itemize subitem}{fg=primary-gold}
  \setbeamercolor{alerted text}{fg=primary-gold}
  \setbeamercolor{block title}{fg=white, bg=primary-blue}
  \setbeamercolor{block body}{bg=light-bg}
  % Semantic block variants: block=definition/concept (blue), exampleblock=
  % worked example (green/positive), alertblock=key takeaway or caution (gold).
  % Keep to <=2 colored boxes per slide across ALL three variants (INV-7).
  \setbeamercolor{block title example}{fg=white, bg=positive}
  \setbeamercolor{block body example}{bg=light-bg}
  \setbeamercolor{block title alerted}{fg=white, bg=primary-gold}
  \setbeamercolor{block body alerted}{bg=light-bg}

  % Minimal footer: course code (left) + page X / N (right), no navigation clutter
  \setbeamertemplate{navigation symbols}{}
  \setbeamertemplate{footline}{%
    \color{neutral}\footnotesize\hspace{0.7em}\@coursecodetext\hfill
    \insertframenumber/\inserttotalframenumber\hspace{0.7em}\vspace{0.3em}}
}{}
\makeatother

% -----------------------------------------------------------------------------
% TikZ setup — libraries the snippet gallery uses
% -----------------------------------------------------------------------------
\usepackage{tikz}
\usetikzlibrary{arrows.meta, positioning, calc, decorations.pathreplacing,
                fit, shapes.geometric, backgrounds}

% Shared TikZ styles — snippets and hand-written diagrams can pick these up
% via `\begin{tikzpicture}[every node/.style={...}]` or by naming specific
% styles (e.g., `\node[dag-node] (X) at (0,0) {X};`).
\tikzset{
  % Node styles for causal DAGs and similar
  dag-node/.style = {
    draw=primary-blue, fill=light-bg, rounded corners=2pt,
    minimum width=1.2cm, minimum height=0.9cm, align=center, font=\small
  },
  decision-node/.style = {
    draw=primary-gold, fill=white, diamond, aspect=1.8,
    minimum width=1.8cm, minimum height=1.2cm, align=center, font=\small,
    inner sep=1pt
  },
  % Edge styles — observed vs counterfactual
  observed-edge/.style = {-{Latex[length=2.5mm]}, thick, draw=jet},
  counterfactual-edge/.style = {-{Latex[length=2.5mm]}, thick, draw=neutral, dashed},
  confound-edge/.style = {-{Latex[length=2.5mm]}, thick, draw=negative, dashed},
  % Data-point styles for plots
  observed-dot/.style = {fill=primary-blue, draw=primary-blue, circle, inner sep=1.2pt},
  counterfactual-dot/.style = {fill=white, draw=neutral, circle, inner sep=1.2pt}
}

% -----------------------------------------------------------------------------
% Convenience macros
% -----------------------------------------------------------------------------
\newcommand{\muted}[1]{\textcolor{neutral}{#1}}
\newcommand{\key}[1]{\textcolor{primary-gold}{\textbf{#1}}}
\newcommand{\good}[1]{\textcolor{positive}{#1}}
\newcommand{\bad}[1]{\textcolor{negative}{#1}}

% Standout frame macro: full-bleed dark background for section transitions.
% Beamer-only (uses \begin{frame}/\setbeamercolor/\usebeamerfont) -- do not
% call from an article-class file (e.g. Notes/<CODE>/*-notes.tex). \muted,
% \key, \good, \bad, and \coursecode above are plain xcolor/gdef and are
% safe to use from either class; verified when Notes/article-class support
% was added.
% Expand: \transitionslide{Your transition title}
\newcommand{\transitionslide}[1]{%
  \begingroup
  \setbeamercolor{background canvas}{bg=primary-blue}
  \begin{frame}[plain]
    \centering
    \vfill
    {\usebeamerfont{title}\color{white}\Large #1\par}
    \vfill
  \end{frame}
  \endgroup
}

% Section divider frame (Beamer-only), an alternative to \transitionslide with a
% darker two-line style: near-black (jet) background, a small white label line
% above a larger gold title, and a thin gold rule along the bottom edge.
% Expand: \sectiondivider{Section 2}{Pointers}
\newcommand{\sectiondivider}[2]{%
  \begingroup
  \setbeamercolor{background canvas}{bg=jet}
  \begin{frame}[plain]
    \vspace*{3.0cm}
    {\color{white}\ttfamily\large #1\par}
    \vspace{0.5cm}
    {\color{primary-gold}\LARGE #2\par}
    \begin{tikzpicture}[remember picture, overlay]
      \draw[primary-gold, line width=2.5pt]
        ($(current page.south west)+(0,0.1)$) -- ($(current page.south east)+(0,0.1)$);
    \end{tikzpicture}
  \end{frame}
  \endgroup
}

% -----------------------------------------------------------------------------
% Channel watermark — "Concepts That Click" (YouTube).
%
% Article-class files (CompetitiveExam Q/A sets, Notes, Assignments) get a
% light diagonal draft watermark on every page plus a centered footer naming
% the channel. Beamer decks get a subtle TikZ background watermark instead
% (the footline already carries the course code). Centralized here so every
% MCQ PDF is branded without per-file edits.
% -----------------------------------------------------------------------------
\makeatletter
\@ifclassloaded{article}{%
  \usepackage{draftwatermark}
  \SetWatermarkText{Concepts That Click}
  \SetWatermarkScale{1}
  \SetWatermarkFontSize{2cm}
  \SetWatermarkLightness{0.86}
  \SetWatermarkAngle{45}
  \usepackage{fancyhdr}
  \pagestyle{fancy}
  \fancyhf{}
  \fancyfoot[C]{\footnotesize\textcolor{neutral}{Concepts That Click~|~YouTube: \href{https://www.youtube.com/@concepts.thatclick}{@concepts.thatclick}}}
  \renewcommand{\headrulewidth}{0pt}
  \renewcommand{\footrulewidth}{0pt}
  \fancypagestyle{plain}{%
    \fancyhf{}%
    \fancyfoot[C]{\footnotesize\textcolor{neutral}{Concepts That Click~|~YouTube: \href{https://www.youtube.com/@concepts.thatclick}{@concepts.thatclick}}}%
    \renewcommand{\headrulewidth}{0pt}%
    \renewcommand{\footrulewidth}{0pt}%
  }
}{}
\@ifclassloaded{beamer}{%
  \setbeamertemplate{background}{%
    \begin{tikzpicture}[remember picture, overlay]
      \node[rotate=45, opacity=0.055, align=center]
        at (current page.center)
        {\Huge\textcolor{neutral}{Concepts That Click}};
    \end{tikzpicture}%
  }
}{}
\makeatother 
\coursecode{JKSSB}
\renewcommand{\thesection}{1.\arabic{section}}
\newtheorem{example}{Example}[section]
\newenvironment{definitionbox}[1]{%
  \par\noindent\textbf{Definition (#1).}\ \itshape
}{\par\vspace{0.3em}}
\title{What Is a Computer? --- Detailed Lecture Notes}
\author{JKSSB Computer Awareness}
\date{\today}

\begin{document}
\maketitle
\noindent\textbf{Video lesson:}
\href{https://youtu.be/EdRZK57oqpo}{What is a Computer?}

\section{A computer is a data-processing system}
Start with an office task. A clerk enters marks into a spreadsheet, a formula
adds the values and calculates an average, the result appears on screen, and
the workbook can be saved. A computer is not merely the screen or the box on
the desk. The task depends on data, instructions, processing, an output path,
and possibly storage.

\begin{definitionbox}{Computer}
An electronic, programmable system that accepts data, processes it under
instructions, produces information, and can store data or results.
\end{definitionbox}

Each part of the definition matters. \emph{Electronic} describes the
technology used by the system. \emph{Programmable} means that instructions
specify the operations to perform. The system accepts inputs, applies those
instructions, and makes results available. It can also retain data or results
for later use.

\begin{definitionbox}{Data}
Raw facts or symbols supplied to a system before they have been organized for
the particular question being asked.
\end{definitionbox}

\begin{definitionbox}{Information}
Data that have been organized or processed so that they are useful for a
particular purpose.
\end{definitionbox}

\begin{definitionbox}{Instruction}
A rule that tells the computer what operation to perform or what action to
take.
\end{definitionbox}

\begin{definitionbox}{Program}
An organized set of instructions that directs a computer through a task.
\end{definitionbox}

Data and information are defined by their role in a task, not by whether the
values are numbers or text. Individual marks are data when the question is
``What are the three scores?'' Their average is information when the question
is ``What was the student's average?'' The average can itself become input
data for a later task, so the same value may play different roles in
different contexts.

\begin{example}
A clerk enters $68$, $74$, and $81$. These three marks are the input data.
Adding them gives $223$; dividing by three gives
$223/3=74.333\ldots$, or approximately $74.33$ to two decimal places. The
total and average are processed results. They are useful information if the
task is to summarize performance. Saving the workbook retains the records.
\end{example}

\begin{example}
Suppose the same three marks are each out of $100$. The total possible score
is therefore $300$. The percentage is
\[
  \frac{68+74+81}{300}\times 100
  =\frac{223}{300}\times 100
  =74.333\ldots\%.
\]
Rounded to two decimal places, the percentage is $74.33\%$. The total,
maximum possible score, formula, and rounding rule all matter; the computer
cannot infer a missing maximum score from the marks alone.
\end{example}

\section{Input, processing, output, and storage}
\begin{definitionbox}{Input}
Data and instructions that enter a computer system through a device, file,
network, or software interface.
\end{definitionbox}

\begin{definitionbox}{Processing}
Operations carried out on input data according to instructions, usually by
the processor and program involved in the task.
\end{definitionbox}

\begin{definitionbox}{Output}
The result made available by a computer system, for example on a display, in
a printed document, as sound, or through another system.
\end{definitionbox}

\begin{definitionbox}{Storage}
The retention of data, instructions, or results for current or later use.
\end{definitionbox}

The basic three-stage model is IPO: Input, Processing, Output. Many
introductory diagrams add Storage, giving IPOS. Storage is a distinct
function, not another name for processing. The stages describe roles in a
task; they do not require every computer job to follow a single rigid,
one-direction path. A saved result can be read back as input to a later
operation.

\begin{center}
\begin{tabular}{p{0.17\textwidth}p{0.34\textwidth}p{0.37\textwidth}}
\toprule
\textbf{Stage} & \textbf{Question to ask} & \textbf{Marks-spreadsheet example}\\
\midrule
Input & What data or instructions enter? & The three marks and the average formula\\
Processing & What operation changes or organizes the data? & Addition, division, and percentage calculation\\
Output & What result is presented, and how? & The total or percentage shown on screen or printed\\
Storage & What is retained for later use? & The saved workbook and its records\\
\bottomrule
\end{tabular}
\end{center}

The physical item does not determine the stage by itself. A monitor is
normally an output device, but the classification question concerns the
monitor's role in the described task. A saved file is stored data; opening
that same file later makes its contents input to the new task. Likewise,
information describes the usefulness of processed data, whereas output
describes how the result is made available. A screen can present useful
information, but the terms answer different questions.

\begin{example}
In the marks spreadsheet, typing or importing the scores is input. The
formula performs processing. The displayed total and average are output.
Saving the workbook is storage. If the saved workbook is opened tomorrow,
reading it is input to that later session.
\end{example}

\section{Characteristics and limitations}
\begin{definitionbox}{Speed}
The ability to carry out many operations in a short time.
\end{definitionbox}

\begin{definitionbox}{Accuracy}
The ability to produce consistent, precise results when the data, instructions,
and operation are correct.
\end{definitionbox}

\begin{definitionbox}{Diligence}
The ability to repeat work without tiring in the human sense, provided the
system continues to receive power and function correctly.
\end{definitionbox}

\begin{definitionbox}{Versatility}
The ability to perform different kinds of tasks by running different
programs.
\end{definitionbox}

\begin{definitionbox}{Automation}
The ability to continue a programmed operation after it has been started and
the required inputs or conditions have been supplied.
\end{definitionbox}

\begin{definitionbox}{Storage as a characteristic}
The ability to retain data, instructions, and results for later access.
\end{definitionbox}

These characteristics describe capabilities, not a guarantee that every
result is correct or useful. Speed concerns how quickly operations are
performed; accuracy concerns whether the operation is correctly carried out
on the supplied data. Storage concerns retaining material; it does not
validate that material. Versatility comes from using different programs,
while automation means a started procedure can continue according to its
instructions and conditions.

Speed and accuracy must not be confused with judgement. A computer can
calculate quickly and repeat the same operation consistently, but it cannot
infer that a deliberately or accidentally incorrect mark is false merely
because the value is stored in a spreadsheet. Accuracy is conditional on the
quality of data and instructions.

\begin{definitionbox}{GIGO}
``Garbage In, Garbage Out'': unsuitable or incorrect input can lead to
unsuitable or incorrect output, even when the computer follows its
instructions.
\end{definitionbox}

\begin{example}
Suppose the correct marks are $68$, $74$, and $81$, but a user enters $18$
instead of $81$. The formula may correctly compute
$68+74+18=160$ and $160/3\approx 53.33$. The arithmetic is consistent with
the entered values, but the output is not a valid summary of the actual
marks. The error entered at the input stage was not corrected by processing.
\end{example}

The example illustrates the limitation often summarized by GIGO. If the
input is wrong, a correct formula can still produce an incorrect result. If
the input is correct but the formula is wrong, the result can also be wrong.
The machine may execute either set of instructions consistently; it does not
automatically decide that the source data are truthful, complete, or fair.
Computers do not possess human judgement or common-sense understanding by
default, so a person must interpret output in context and decide what action
it justifies.

\section{Applications and role classification}
The same IPO/IPOS model applies across fields. What changes is the type of
data and the useful result; the roles of input, processing, output, and
storage remain identifiable.

\begin{center}
\begin{tabular}{p{0.15\textwidth}p{0.24\textwidth}p{0.27\textwidth}p{0.22\textwidth}}
\toprule
\textbf{Field} & \textbf{Input} & \textbf{Processing / output} & \textbf{Storage}\\
\midrule
Office & Text, figures, and instructions & Create a document or calculate a total & Document or workbook file\\
Banking & Account and transaction details & Validate and process a transaction; show a status & Account and transaction records\\
Education & Learner responses and assessment rules & Score work and present results or learning resources & Assessment and learner records\\
Healthcare & Patient details and appointment requests & Organize records, schedules, or diagnostic support & Patient and appointment records\\
Governance & Online form fields and supporting details & Process a request and provide a citizen service & Submitted forms and service records\\
\bottomrule
\end{tabular}
\end{center}

For exam questions, classify the role within the described task. A keyboard
is an input device; the formula is processing; a monitor presents output; a
saved file is stored data. A physical object may support more than one
stage, but the question usually asks for the role it is playing in that
particular scenario.

\begin{example}
An online expense form accepts entered amounts, checks and sums them using
the programmed rules, displays the total, and stores the submitted record.
If one amount is mistyped, the calculation may still be executed correctly
on the wrong data. The user should correct the input before relying on the
total.
\end{example}

When a question gives a complete workflow, mark the stage of each action
rather than assigning one label to the entire system. Entering a value is
input; checking or summing it is processing; showing a result is output;
keeping the transaction record is storage. If the stored record is opened
and processed again, its role changes to input for that later operation.

\subsection*{Exercises}
\begin{enumerate}[label=\arabic*.]
  \item A user enters four prices, applies a tax formula, prints the invoice,
    and saves the file. Identify the input, processing, output, and storage.
  \item Explain why a fast, consistent computer can still produce an
    incorrect result.
  \item Three temperature readings are $18$, $20$, and $22$ degrees. For a
    task asking for their average, identify the data and the resulting
    information.
  \item Give one example from banking or governance and trace it through
    IPOS.
\end{enumerate}

\subsection*{Solutions}
\begingroup\small
\begin{enumerate}[label=\arabic*.]
  \item The four prices are input; applying the tax formula is processing;
    the printed invoice is output; saving the invoice file is storage.
  \item The computer follows its instructions and operates on the supplied
    data. Incorrect input or a wrong formula can therefore produce an
    incorrect result consistently; speed does not establish truth.
  \item The three readings are input data. Their average is
    $(18+20+22)/3=20$ degrees and is information for the average-reading
    task.
  \item Example: an online fee form accepts the user's details and amount
    (input), calculates the fee (processing), shows a receipt (output), and
    saves the transaction record (storage).
\end{enumerate}
\endgroup
\end{document}
